% ============================================================
% ch05_results.tex
% Demo Flow, Advantages, Limitations, Viva Preparation
% ============================================================

\chapter{Demonstration Flow}

\section{Recommended Viva Demo Sequence}

The following sequence demonstrates all ADBMS concepts within 10--15 minutes:

\begin{table}[H]
\centering
\caption{Demo Sequence for Viva}
\begin{tabularx}{\textwidth}{clX}
\toprule
\textbf{Step} & \textbf{Menu Option} & \textbf{Concept Shown} \\
\midrule
1 & 7 (Show Balances) & Starting PostgreSQL state \\
2 & 6 (Show WAL Logs) & Empty or existing MongoDB logs \\
3 & 1 (Deposit) & WAL written → PG updated → COMMITTED \\
4 & 6 (Show WAL Logs) & See the COMMITTED entry \\
5 & 2 (Withdraw, large amount) & Insufficient funds → ROLLBACK → ROLLED\_BACK \\
6 & 4 (Simulate Crash) & UNCOMMITTED log created in MongoDB \\
7 & 6 (Show WAL Logs) & Point to the UNCOMMITTED entry \\
8 & 7 (Show Balances) & Verify PG balance unchanged (auto-rollback) \\
9 & 5 (Run Recovery) & Analysis → UNDO → REDO phases \\
10 & 6 (Show WAL Logs) & Entry now ROLLED\_BACK \\
11 & 7 (Show Balances) & Consistent state confirmed \\
\bottomrule
\end{tabularx}
\end{table}

\section{Sample Terminal Output}

\subsection{Normal Deposit (Option 1)}

\begin{lstlisting}[style=jsonstyle, caption={Deposit terminal output}]
  [Session] sess_a1b2c3

  ====================================================
  DEPOSIT | ACC-00001 | +1500.00
  ====================================================
  Old balance: 10000.00
    [WAL] Log #1 -> UNCOMMITTED  (DEPOSIT, account=1, 10000.00->11500.00)
    [PG]  COMMIT -- Transaction_Record id=13
    [WAL] Log #1 -> COMMITTED
  New balance: 11500.00
  v Deposit successful.
\end{lstlisting}

\subsection{Rollback --- Insufficient Funds (Option 2)}

\begin{lstlisting}[style=jsonstyle, caption={Rollback terminal output}]
  ====================================================
  WITHDRAWAL | ACC-00005 | -50000.00
  ====================================================
  Old balance: 8000.00
    [WAL] Log #2 -> UNCOMMITTED  (WITHDRAWAL, account=5, 8000.00->-42000.00)
  [ERROR] Insufficient funds. Balance: 8000.00, Requested: 50000.00
  [PG] ROLLBACK -- balance unchanged.
    [WAL] Log #2 -> ROLLED_BACK
\end{lstlisting}

\subsection{Crash Simulation (Option 4)}

\begin{lstlisting}[style=jsonstyle, caption={Crash simulation terminal output}]
  ====================================================
  CRASH SIMULATION -- WITHDRAWAL
  Account : ACC-00001  |  Amount: 3000.00
  ====================================================
  Current balance: 11500.00
    [WAL] Log #3 -> UNCOMMITTED  (WITHDRAWAL, account=1, ...)
  [PG]  Balance changed: 11500.00 -> 8500.00
  [PG]  (Change applied but NOT yet committed)

  *** SYSTEM CRASH -- Exception raised! ***
  *** COMMIT was never executed.        ***

  [PG]  Connection closed -> PG auto-rollback
        (PG balance restored to 11500.00)
  [WAL] Log #3 still UNCOMMITTED in MongoDB
  -> Run Recovery Manager (option 5) to fix the database.
\end{lstlisting}

\subsection{Recovery Manager (Option 5)}

\begin{lstlisting}[style=jsonstyle, caption={Recovery Manager terminal output}]
  ####################################################
  RECOVERY MANAGER -- Starting Recovery
  ####################################################

  -- Phase 1: ANALYSIS --
  Scanning MongoDB WAL logs...
  Found 1 UNCOMMITTED log(s) -> need UNDO
  Found 5 COMMITTED   log(s) -> check REDO

  Crash victims (will be UNDONE):
    log_id=3  account_id=1  type=WITHDRAWAL  old=11500.00  new=8500.00

  -- Phase 2: UNDO --
  Restoring 1 account(s) to pre-crash state...

  UNDO log_id=3  account_id=1
    PG balance now : 11500.00
    Restoring to   : 11500.00  <- from WAL log
    Already at old_balance. UNDO already done. Skipping.

  -- Phase 3: REDO --
  v All committed transactions consistent. No REDO needed.

  ####################################################
  RECOVERY COMPLETE
  UNDO applied : 0   (PG already consistent -- auto-rollback handled it)
  REDO applied : 0
  PostgreSQL is now in a consistent state.
  ####################################################
\end{lstlisting}

% ────────────────────────────────────────────────────────────
\chapter{Advantages and Limitations}

\section{Advantages}

\begin{enumerate}
  \item \textbf{Educational Clarity}: Every step of WAL and recovery is
        visible in the terminal output. Students can trace exactly what
        happens to both databases at each transaction stage.

  \item \textbf{Polyglot Persistence}: Demonstrates a hybrid architecture
        where two databases complement each other --- a real-world pattern
        used in modern applications.

  \item \textbf{Practical Implementation}: The entire system runs on a
        single laptop using freely available software (PostgreSQL, MongoDB,
        Python).

  \item \textbf{Observable Recovery}: Unlike built-in DBMS recovery which
        is invisible to users, this system makes recovery explicit and
        observable through MongoDB queries and terminal output.

  \item \textbf{Idempotent Recovery}: The recovery manager can be run
        multiple times safely, demonstrating a key property of robust
        recovery systems.

  \item \textbf{Modular Design}: Each Python module is focused on one
        responsibility, making the code easy to understand, extend, and
        explain during viva.
\end{enumerate}

\section{Limitations}

\begin{enumerate}
  \item \textbf{Single-User Only}: The system does not handle concurrent
        transactions. Running two transactions simultaneously could
        produce incorrect results.

  \item \textbf{No Checkpoint Support}: A full ARIES implementation uses
        checkpoints to limit how far back in the log recovery must scan.
        Our system always scans all logs.

  \item \textbf{No LSN Chain}: Log Sequence Numbers link log records for
        ordered undo/redo. Our system uses timestamps and log\_id ordering,
        which is less precise.

  \item \textbf{Simplified Crash Model}: Our crash simulation uses a
        Python exception. A real crash (power loss, OS kill) has more
        complex storage-level effects.

  \item \textbf{No Page-Level Recovery}: Our system operates at the
        account/row level. Real DBMS recovery operates at the page level
        for finer-grained control.

  \item \textbf{MongoDB Not Crash-Safe by Default}: If MongoDB itself
        crashes before the log insert completes, the log is also lost.
        A production system would require MongoDB with journaling enabled
        and write concern set to ``majority''.
\end{enumerate}

% ────────────────────────────────────────────────────────────
\chapter{Viva Preparation}

\section{Expected Teacher Questions and Answers}

\subsection{Q1: What is Write-Ahead Logging?}

\textbf{Answer:} Write-Ahead Logging is the rule that a log record must be
written to stable storage \textit{before} the data change it describes is
applied. In our system, we insert a document into MongoDB's \texttt{wal\_logs}
collection before running the \texttt{UPDATE Account SET balance = ...}
statement in PostgreSQL. This guarantees that even if PostgreSQL crashes
mid-update, the log exists and we know what to undo.

\subsection{Q2: What is the difference between UNDO and REDO?}

\textbf{Answer:}
\begin{itemize}
  \item \textbf{UNDO}: Reverses an \textit{uncommitted} transaction by
        restoring \texttt{old\_balance} (read from the WAL log). Used for
        crash victims that never reached COMMIT.
  \item \textbf{REDO}: Re-applies a \textit{committed} transaction by
        applying \texttt{new\_balance} (from the WAL log). Used when a
        committed change was lost before being flushed to PostgreSQL disk.
\end{itemize}

\subsection{Q3: How do you identify uncommitted transactions after a crash?}

\textbf{Answer:} We query MongoDB for all documents where
\texttt{status = "UNCOMMITTED"}. These are transactions that wrote a WAL
log but never had their status updated to \texttt{COMMITTED} or
\texttt{ROLLED\_BACK} --- the crash happened between those steps.

\subsection{Q4: Why use MongoDB for WAL logs instead of PostgreSQL?}

\textbf{Answer:} MongoDB is used because:
(1) Log inserts are fast and lightweight without transaction overhead.
(2) MongoDB is independent of the PostgreSQL transaction --- a PG rollback
does not affect the MongoDB log.
(3) JSON documents are flexible and human-readable, making logs easy to
inspect and demonstrate.
(4) It demonstrates polyglot persistence --- using each database for what
it does best.

\subsection{Q5: What is ACID? How does your system demonstrate it?}

\textbf{Answer:}
\begin{itemize}
  \item \textbf{Atomicity}: Transfer uses one \texttt{conn.commit()} for
        both account updates. If one fails, both rollback.
  \item \textbf{Consistency}: CHECK constraints and FK constraints in
        PostgreSQL prevent invalid data (e.g., negative balance).
  \item \textbf{Isolation}: PostgreSQL handles concurrent isolation via
        its transaction model (one connection at a time in our demo).
  \item \textbf{Durability}: WAL log is written first, ensuring the
        operation can be recovered even after a crash.
\end{itemize}

\subsection{Q6: How is your recovery similar to ARIES?}

\textbf{Answer:} ARIES has three phases: Analysis, Redo, Undo. Our system
follows the same three phases. The difference is that ARIES uses Log Sequence
Numbers for precise ordering, maintains a dirty page table, and operates at
the page level. Our system uses timestamps and log\_id for ordering, operates
at the account row level, and is simplified for undergraduate demonstration.

\subsection{Q7: What is polyglot persistence?}

\textbf{Answer:} Polyglot persistence means using multiple database
technologies in one application, each chosen for its strengths. In our
system, PostgreSQL handles ACID-compliant relational banking data while
MongoDB handles flexible, fast-insert WAL log documents. Together they form
a more capable system than either alone.

\subsection{Q8: Show me the WAL log in MongoDB right after a crash.}

\textbf{Answer:} Run option 4 from the menu, then in mongosh run:
\begin{lstlisting}[style=jsonstyle]
use("bank_recovery_logs")
db.wal_logs.find({ status: "UNCOMMITTED" })
\end{lstlisting}
This shows the log entry written before the crash, still showing
\texttt{UNCOMMITTED} because the system never reached the COMMIT step.

\section{Key Concepts Quick Reference}

\begin{table}[H]
\centering
\caption{Quick Reference for Viva}
\begin{tabularx}{\textwidth}{lX}
\toprule
\textbf{Concept} & \textbf{One-line explanation} \\
\midrule
WAL & Log first, then update the database \\
UNCOMMITTED & Transaction started but crash happened before COMMIT \\
UNDO & Restore old\_balance to reverse an uncommitted transaction \\
REDO & Re-apply new\_balance for a committed but lost transaction \\
Atomicity & All steps of a transaction succeed or all are reversed \\
Durability & Committed data survives crashes (WAL makes this possible) \\
Polyglot & Two databases used together, each for its strength \\
Idempotency & Recovery gives same result whether run once or many times \\
\bottomrule
\end{tabularx}
\end{table}
